\begin{definitionbox}{Definition ($M$-Tail of a Sequence)}
% ---------------------------------------------------------
% Definition: $M$-Tail of a Sequence
% Label: def:m-tail
% ---------------------------------------------------------
\begin{definition}[$M$-Tail of a Sequence]
\label{def:m-tail}

Let $(x_n) \subseteq \mathbb{R}$ and let $M \in \mathbb{N}$. The
\textbf{$M$-tail} of $(x_n)$ is the sequence
\[
    (x_n)_{n \ge M} \;:=\; \bigl(x_M,\, x_{M+1},\, x_{M+2},\, \ldots\bigr).
\]
\end{definition}
\end{definitionbox}

\begin{remark*}[Standard quantified statement]
\[
\forall (x_n) \subseteq \mathbb{R} \;\forall M \in \mathbb{N}
\;\Bigl(
    (x_n)_{n \ge M} = (x_M,x_{M+1},x_{M+2},\ldots)
\Bigr).
\]
\end{remark*}

\begin{remark*}[Definition predicate reading]
\[
\operatorname{TailSequence}(x_n,M,y_k)
\;:=\;
\forall k \in \mathbb{N}\;\bigl(y_k = x_{M+k-1}\bigr).
\]
\end{remark*}

\begin{remark*}[Interpretation]
The $M$-tail discards the first $M - 1$ terms and retains everything
from index $M$ onward. It is the formal object underlying the phrase
\emph{all but finitely many terms}: to say that all but finitely many
terms of $(x_n)$ satisfy a property $P$ is precisely to say that some
$M$-tail of $(x_n)$ satisfies $P$ termwise. The index $M$ is an offset
into the sequence, not a bound on the terms themselves; two sequences
with identical $M$-tails share all asymptotic properties, differing
only in finitely many initial values which convergence arguments
discard entirely.
\end{remark*}

\begin{dependencies}
\begin{itemize}
  \item \hyperref[def:real-sequence]{Real Sequence}
\end{itemize}
\end{dependencies}



\begin{theorembox}{Theorem (Convergence of a Tail)}
% ---------------------------------------------------------
% Theorem: Convergence of a Tail
% Label: thm:convergence-of-tail
% ---------------------------------------------------------
\begin{theorem}[Convergence of a Tail]
\label{thm:convergence-of-tail}

Let $(x_n) \subseteq \mathbb{R}$ and let $m \in \mathbb{N}$. The
$m$-tail $(x_{m+n})_{n \in \mathbb{N}}$ converges if and only if
$(x_n)$ converges. Moreover, in this case,
\[
    \lim_{n \to \infty} x_{m+n} \;=\; \lim_{n \to \infty} x_n.
\]

\smallskip
\noindent
\hyperref[prf:convergence-of-tail]{\textit{Go to proof.}}
\end{theorem}
\end{theorembox}

\begin{remark*}[Standard quantified statement]
\[
\begin{aligned}
&\forall (x_n) \subseteq \mathbb{R} \;\forall m \in \mathbb{N}\\
&\quad\Bigl(
    \bigl(x_{m+n}\bigr)_{n \in \mathbb{N}} \text{ converges}
    \;\Longleftrightarrow\;
    (x_n) \text{ converges}
\Bigr)\\
&\quad\land\;\Bigl(
    (x_n) \text{ converges}
    \;\Longrightarrow\;
    \lim_{n\to\infty} x_{m+n} = \lim_{n\to\infty} x_n
\Bigr).
\end{aligned}
\]
\end{remark*}

\begin{remark*}[Theorem predicate reading]
\[
\operatorname{TailConvergenceInvariant}(x_n,m)
\;:=\;
\Bigl(\exists L\in\mathbb{R}\,\operatorname{ConvergentSequence}(x_n,L)\Bigr)
\Longleftrightarrow
\Bigl(\exists L\in\mathbb{R}\,\operatorname{ConvergentSequence}(x_{m+n},L)\Bigr).
\]
\end{remark*}

\begin{remark*}[Contrapositive quantified statement]
\[
\begin{aligned}
&\forall (x_n) \subseteq \mathbb{R} \;\forall m \in \mathbb{N}\\
&\quad\Bigl(
    \bigl(x_{m+n}\bigr)_{n \in \mathbb{N}} \text{ diverges}
    \;\Longleftrightarrow\;
    (x_n) \text{ diverges}
\Bigr).
\end{aligned}
\]
\end{remark*}

\begin{remark*}[Contrapositive predicate reading]
\[
\lnot\operatorname{TailConvergenceInvariant}(x_n,m)
\;:=\;
\bigl((x_n)\text{ diverges}\Longleftrightarrow (x_{m+n})_{n\in\mathbb{N}}\text{ diverges}\bigr).
\]
\end{remark*}

\begin{remark*}[Interpretation]
Convergence is a \emph{tail property}: it depends only on the
eventual behavior of the sequence, not on any finite initial
segment. Discarding the first $m$ terms neither creates nor
destroys convergence, and when the limit exists it is preserved
exactly. This theorem formalizes the informal principle that
finitely many terms are irrelevant to asymptotic behavior. It is
the rigorous justification for the common proof move of assuming
without loss of generality that $n$ is large enough for some
auxiliary condition to hold: one simply passes to a suitable
$m$-tail, applies the argument there, and concludes for the full
sequence by this theorem. The contrapositive is equally useful:
to show $(x_n)$ diverges it suffices to show that some $m$-tail
diverges.
\end{remark*}

\begin{dependencies}
\begin{itemize}
  \item \hyperref[def:convergent-sequence]{Convergence of a Sequence}
  \item \hyperref[def:m-tail]{$M$-Tail of a Sequence}
  \item \hyperref[thm:uniqueness-of-limits]{Uniqueness of Limits}
\end{itemize}
\end{dependencies}









\begin{theorembox}{Theorem (Convergence by Domination)}
% ---------------------------------------------------------
% Theorem: Convergence by Domination
% Label: thm:convergence-by-domination
% ---------------------------------------------------------
\begin{theorem}[Convergence by Domination]
\label{thm:convergence-by-domination}

Let $(x_n) \subseteq \mathbb{R}$ and let $L \in \mathbb{R}$. Let
$(a_n)$ be a sequence of positive real numbers with
$\displaystyle\lim_{n \to \infty} a_n = 0$. If there exist a constant
$c > 0$ and an index $m \in \mathbb{N}$ such that
\[
    |x_n - L| \;\le\; c\, a_n \qquad \forall n \ge m,
\]
then $\displaystyle\lim_{n \to \infty} x_n = L$.

\smallskip
\noindent
\hyperref[prf:convergence-by-domination]{\textit{Go to proof.}}
\end{theorem}
\end{theorembox}

\begin{remark*}[Standard quantified statement]
\[
\begin{aligned}
&\forall (x_n) \subseteq \mathbb{R} \;\forall L \in \mathbb{R}
 \;\forall (a_n) \subseteq (0, \infty)\\
&\quad\Bigl(
    \Bigl(
        \lim_{n \to \infty} a_n = 0
        \;\land\;
        \exists c > 0 \;\exists m \in \mathbb{N}
        \;\forall n \ge m \;\bigl(|x_n - L| \le c\, a_n\bigr)
    \Bigr)
    \;\Longrightarrow\;
    \lim_{n \to \infty} x_n = L
\Bigr).
\end{aligned}
\]
\end{remark*}

\begin{remark*}[Contrapositive quantified statement]
\[
\begin{aligned}
&\forall (x_n) \subseteq \mathbb{R} \;\forall L \in \mathbb{R}
 \;\forall (a_n) \subseteq (0,\infty)\\
&\quad\Bigl(
    \lim_{n \to \infty} x_n \neq L
    \;\Longrightarrow\;
    \lim_{n \to \infty} a_n \neq 0
    \;\lor\;
    \forall c > 0 \;\forall m \in \mathbb{N}
    \;\exists n \ge m \;\bigl(|x_n - L| > c\, a_n\bigr)
\Bigr).
\end{aligned}
\]
\end{remark*}

\begin{remark*}[Contrapositive predicate reading]
\[
\operatorname{NotDominatedByNullError}(x_n,L,a_n)
\;:=\;
\lim_{n\to\infty}a_n\neq 0
\lor
\forall c>0\,\forall m\in\mathbb{N}\,\exists n\ge m
\bigl(|x_n-L|>c\,a_n\bigr).
\]
\end{remark*}

\begin{remark*}[Interpretation]
The theorem is a \emph{domination principle}: it is not necessary to
verify the $\varepsilon$-$N$ condition directly for $(x_n)$. It
suffices to find a null sequence $(a_n)$ -- one converging to zero --
that dominates the error $|x_n - L|$ on some tail, up to a fixed
scalar $c$. The scalar $c$ is irrelevant to the conclusion: since
$a_n \to 0$, the product $c\, a_n \to 0$ regardless of the size of
$c$, and any $\varepsilon > 0$ is eventually beaten by $c\, a_n$.
This theorem is the rigorous engine behind the common proof pattern
of bounding $|x_n - L|$ by a simpler expression and observing that
the simpler expression tends to zero. The canonical instance is
$a_n = 1/n$, where domination by $c/n$ is sufficient to conclude
convergence to $L$.
\end{remark*}

\begin{dependencies}
\begin{itemize}
  \item \hyperref[def:convergent-sequence]{Convergence of a Sequence}
  \item \hyperref[def:m-tail]{$M$-Tail of a Sequence}
\end{itemize}
\end{dependencies}



\begin{theorembox}{Theorem (Ratio Limit Less Than One Implies Null)}
% ---------------------------------------------------------
% Theorem: Ratio Limit Less Than One Implies Null
% Label: thm:ratio-limit-less-than-one-implies-null
% ---------------------------------------------------------
\begin{theorem}[Ratio Limit Less Than One Implies Null]
\label{thm:ratio-limit-less-than-one-implies-null}

Let $(x_n)$ be a sequence of positive real numbers. Suppose that the
limit
\[
    L := \lim_{n\to\infty}\frac{x_{n+1}}{x_n}
\]
exists. If $L<1$, then $(x_n)$ converges and
\[
    \lim_{n\to\infty} x_n = 0.
\]

\smallskip
\noindent
\hyperref[prf:ratio-limit-less-than-one-implies-null]{\textit{Go to proof.}}
\end{theorem}
\end{theorembox}

\begin{remark*}[Standard quantified statement]
\[
\begin{aligned}
&\forall (x_n)\subseteq (0,\infty)\;\forall L\in\mathbb{R}\\
&\quad\Bigl(
    \lim_{n\to\infty}\frac{x_{n+1}}{x_n}=L
    \;\land\;
    L<1
\Bigr)
\Longrightarrow
\lim_{n\to\infty}x_n=0.
\end{aligned}
\]
\end{remark*}

\begin{remark*}[Theorem predicate reading]
\[
\bigl(
\operatorname{ConvergentSequence}\bigl(x_{n+1}/x_n,L\bigr)
\land
L<1
\bigr)
\Longrightarrow
\operatorname{NullSequence}(x_n).
\]
\end{remark*}

\begin{remark*}[Negated quantified statement]
\[
\begin{aligned}
&\exists (x_n)\subseteq (0,\infty)\;\exists L\in\mathbb{R}\\
&\quad\Bigl(
    \lim_{n\to\infty}\frac{x_{n+1}}{x_n}=L
    \;\land\;
    L<1
    \;\land\;
    \lim_{n\to\infty}x_n\neq 0
\Bigr).
\end{aligned}
\]
\end{remark*}

\begin{remark*}[Negation predicate reading]
\[
\operatorname{ConvergentSequence}\bigl(x_{n+1}/x_n,L\bigr)
\land
L<1
\land
\operatorname{NotConvergentTo}(x_n,0).
\]
\end{remark*}

\begin{remark*}[Failure modes]
A failure would be a positive sequence whose consecutive ratios settle below
one in the limit but whose terms do not tend to zero. The theorem says this
cannot happen: once the ratio limit is less than one, some tail has ratios
bounded by a fixed number $r<1$, forcing geometric decay.
\end{remark*}

\begin{remark*}[Failure mode decomposition]
\[
\underbrace{(x_n)\subseteq(0,\infty)}_{\text{positive terms}}
\quad
\underbrace{\lim_{n\to\infty}\frac{x_{n+1}}{x_n}=L<1}_{\text{eventual ratio gap}}
\quad
\underbrace{\lim_{n\to\infty}x_n\neq 0}_{\text{fails to be null}}.
\]
The positive-term hypothesis makes each ratio meaningful. The strict gap
$L<1$ gives a tail on which $x_{n+1}\le r x_n$ for some $r<1$, and repeated
iteration then gives domination by a null geometric sequence.
\end{remark*}

\begin{remark*}[Interpretation]
This is the sequence form of the ratio test. The theorem does not merely say
that the terms decrease eventually; it says they decrease at an eventually
geometric rate. The strict inequality $L<1$ is essential. When the ratio
limit equals $1$, the conclusion may hold, as with $1/n$, or fail, as with
the constant sequence $1$.
\end{remark*}

\begin{dependencies}
\begin{itemize}
  \item \hyperref[def:convergent-sequence]{Convergence of a Sequence}
  \item \hyperref[def:null-sequence]{Null Sequence}
  \item \hyperref[def:m-tail]{$M$-Tail of a Sequence}
  \item \hyperref[thm:convergence-by-domination]{Convergence by Domination}
\end{itemize}
\end{dependencies}
